\chapter{Discussion and Policy Implications}
\label{ch:discussion}

\section{Why Headline CPI Pass-Through Is Modest}

The estimated cumulative positive pass-through of 0.021 means that a 10 per cent increase in rupee-denominated oil prices raises monthly headline CPI inflation by about 0.21 percentage points. This is an economically meaningful effect, but it is smaller than what might be expected given India's heavy oil import dependence.

There are several reasons for this. Headline CPI in India assigns about 46 per cent weight to food items, which are driven primarily by domestic agricultural supply conditions, monsoon outcomes, and government procurement policies. Only about 7 per cent of the CPI basket falls directly under fuel and light. Oil prices affect headline CPI through several indirect channels (transport costs embedded in food distribution, input costs for manufactured goods, and the relative price of oil-intensive services), but these indirect channels are noisy and operate with variable lags.

The result is that oil price shocks explain only a modest fraction of monthly CPI variation. This does not mean oil is unimportant; it means that headline CPI is the wrong place to look for a large, sharp signal. The Fuel and Light sub-index analysis in the appendix confirms this interpretation, showing a pass-through estimate about three times larger than headline CPI.

\section{The Exchange Rate Channel}

The Brent-plus-EXR robustness specification provides a sharper picture of oil pass-through in India. When the world oil price in dollars is separated from the rupee exchange rate, the positive oil pass-through estimate increases to 0.028 and becomes marginally significant at 10 per cent. The exchange rate coefficient is significant at 5 per cent ($p = 0.029$), indicating that a 1 per cent depreciation of the rupee against the dollar raises monthly CPI inflation by about 0.04 percentage points, holding oil prices constant.

This finding matters for two reasons. First, it confirms that part of the inflationary impact attributed to oil price shocks in the primary specification is actually coming through the exchange rate. India's heavy oil import bill puts persistent downward pressure on the rupee, creating a secondary inflationary channel that operates independently of the direct cost effect of oil. Second, it implies that the RBI's exchange rate management has important implications for inflation outcomes beyond the standard monetary policy channel.

\section{Deregulation and the Pass-Through Regime}

The diesel deregulation dummy in the main model is negative and highly significant ($-0.339$, $p < 0.001$), indicating that average monthly CPI inflation fell after October 2014. This likely reflects the combined effect of deregulation itself (which ended the period of government-subsidised fuel prices and removed one source of fiscal expansion), the global oil price decline of 2014--16, and the RBI's adoption of flexible inflation targeting in February 2015.

The sub-sample comparison shows slightly higher positive pass-through in the pre-2014 period (0.017 vs.\ 0.010), which seems counterintuitive given that deregulation should have made domestic prices more responsive to international oil prices. However, the pre-2014 period includes the 2007--08 oil price spike, when the government temporarily raised fuel prices sharply despite the regulated pricing framework. The post-2014 period, by contrast, includes several years of low and stable oil prices (2015--17) during which pass-through was difficult to estimate precisely regardless of the pricing regime.

\section{Asymmetry: What Can and Cannot Be Said}

The Wald test for asymmetry fails to reject the null of symmetric pass-through at the 5 per cent level ($p = 0.241$) across all specifications. This needs to be interpreted carefully.

The failure to reject does not prove that pass-through is symmetric. It means that the data are not informative enough to distinguish the two pass-through coefficients at the conventional significance level. Given the noise in headline CPI, the limited sample size, and the fact that the underlying asymmetry (if it exists) may be modest, this is not a surprising result. In fact, Kilian and Vigfusson (2011) have argued that asymmetry is inherently difficult to detect with standard time-series methods, because the test has low power when the true asymmetry is small relative to the noise.

What can be said with confidence is that the point estimates are asymmetric across every specification examined: the primary model, the Brent-plus-EXR model, both sub-samples, the winsorised specification, and the rolling window estimates. The direction is always the same, with positive oil shocks having a larger cumulative effect than negative shocks. This consistent pattern across specifications is itself evidence worth noting, even if no single test crosses the significance threshold.

\section{Policy Implications}

For Indian monetary policy, the results suggest that the RBI is right to monitor oil prices as an inflation risk factor, but the short-run pass-through to headline CPI is not large enough to warrant aggressive monetary policy responses to individual oil price shocks. The estimated effect of a 10 per cent oil price increase is about 0.21 percentage points on monthly inflation. Annualised, this would amount to roughly 0.2 percentage points on annual headline CPI, which is within the normal range of monthly CPI volatility.

The more actionable policy implication concerns the exchange rate. Since rupee depreciation independently raises CPI inflation, episodes of oil price increases that coincide with rupee weakness deserve closer attention than oil price increases during periods of rupee stability. The 2022 experience, when Brent spiked above 100 USD per barrel while the rupee weakened past 80 per dollar, is a case in point.

For fiscal policy, the results indirectly support the case for completing the fuel price deregulation process. The persistence of administered pricing for cooking gas and periodic interventions through excise duty adjustments mean that the pass-through mechanism is still partly dampened by government action. A fully deregulated fuel pricing regime would make the pass-through more transparent and possibly more symmetric, as Pradeep (2022) has argued.
